\appendixchapter{Algèbre multilinéaire}

\section{Produit tensoriel}

\begin{defi}[Produit tensoriel]
Soient $E$,$F$ des $\R$-espaces vectoriels de dimension finie, on appelle produit tensoriel de $E$ et $F$ est on note $E\otimes F$ l'unique (à isomorphisme près) espace vectoriel muni d'une application bilinéaire $(x,y) \in E\times F \mapsto x\otimes y \in E\otimes F$ vérifiant la propriété universelle: \\
Pour toute application bilinéaire $f: E\times F \to G$, il existe une unique application linéaire $g: E\otimes F \to G$ telle que le diagramme suivant commute
$$\begin{tikzcd}
E \times F \arrow[r, "f"]  \arrow[d, ".\otimes ."] & G \\
E\otimes F \arrow[ru, "g"']& 
\end{tikzcd}$$
On appelle l'élément $x\otimes y$ de $E\times F$ un tenseur pur. \\

De plus, si $(e_\mu)_\mu$ est une base de $E$ et $(f_\nu)_\nu$ une base de $F$ alors $(e_\mu\otimes f_\nu)_{\mu,\nu}$ est une base de $E\otimes F$, ainsi $\dim E\otimes F = \dim E \dim F$.
\end{defi}
\begin{preu}
Considérons $\R^{(E\times F)}$ le $\R$-espace vectoriel des suites presques nulles indexées par $E\times F$ dont on note $(e_{x,y})_{(x,y)\in E\times F}$ la base canonique. Soit $\mathcal{R}$ le sous espace vectoriel engendré par les éléments 
$$e_{x+x', y}-e_{x,y}-e_{x',y}$$
$$e_{x, y+y'}-e_{x,y}-e_{x, y'}$$
$$e_{\lambda x, y}-\lambda e_{x, y}$$
$$e_{x, \lambda y}-\lambda e_{x, y}$$
pour $x\in E$, $y\in F$ et $\lambda \in \R$. \\

Alors on note $H = \R^{(E\times F)}/\mathcal{R}$ l'espace vectoriel quotient. De plus l'application $.\otimes . : E\times F \to H$ qui envoie $(x,y)$ sur la classe d'équivalence de $e_{x,y}$ dans $H$ est bilinéaire par définition de $\mathcal{R}$. \\

Soit alors $f:E\times F \to G$ une application bilinéaire. Soit alors l' unique application linéaire $u: \R^{(E\times F)} \to G, e_{x,y} \mapsto f(x,y)$ a un noyau qui contient $\mathcal{R}$ par bilinéarité de $f$, donc l'application $u$ induit une unique application linéaire $g : H\to G$ tel que le diagramme commute 
$$\begin{tikzcd}
E \times F \arrow[r, "f"]  \arrow[d, ".\otimes ."] & G \\
H \arrow[ru, "g"']& 
\end{tikzcd}$$
Ainsi $H$ vérifie la propriété universelle du produit tensoriel donc est unique à isomorphisme près.
\end{preu}

\begin{pro}
Soient $E$, $F$, $G$ et $H$ des $\R$-espaces vectoriels de dimension finie. Alors le produit tensoriel vérifie les propriétés suivantes:
\be
	\item (fonctorialité): Soient $f: E\to G$ et $g: F\to H$ des applications linéaires, alors il existe une unique application linéaire $f\otimes g: E\otimes F \to G \otimes H$ telle que 
	$$f\otimes g(x\otimes y) = f(x) \otimes g(y)$$
	\item $$E \otimes \R \simeq E$$
	via l'isomorphisme qui envoie $e\otimes 1$ sur $e$.
	\item (commutativité): $$E\otimes F \simeq F \otimes E$$
	via l'isomorphisme qui envoie $e\otimes f$ sur $f\otimes e$.
	\item (associativité): $$(E\otimes F)\otimes G \simeq E \otimes (F\otimes G)$$
	via l'isomorphisme qui envoie $(e\otimes f)\otimes g$ sur $e\otimes (f\otimes g)$.
	\item (distributivité): $$E \otimes (F \oplus G) \simeq (E \otimes F) \oplus (E \otimes G)$$
	via l'isomorphisme qui envoie $e\otimes (f\oplus g)$ sur $(e\otimes f) \oplus (e \otimes g)$.
	\item $$E^* \otimes F^* \simeq (E\otimes F)^*$$
	via l'isomorphisme qui envoie $u\otimes v$ sur $x\otimes y \mapsto u(x)v(y)$.
	\item $$E^*\otimes F \simeq \mathscr{L}(E, F)$$
	via l'isomorphisme qui envoie $u\otimes x$ sur $x \mapsto u(x)y$
\ee
\end{pro}

\begin{rem}
Ces isomorphismes permettent d'identifier $E^*\otimes F^*$ avec les formes bilinéaires sur $E \times F$.
\end{rem}

\section{Algèbre extérieure}

\begin{defi}[Algèbre tensorielle]
Soit $V$ un $\R$-espace vectoriel. Alors l'ensemble $\oplus_{k\geq 0} \otimes^k V$ munie de l'application 
$$\begin{array}{ccccc}
	\otimes & : & \otimes^p V\times \otimes^q V & \to & \otimes^{p+q} V \\
		& & (u_1\otimes ...\otimes u_p, v_1 \otimes ... \otimes v_q) & \mapsto & u_1 \otimes ... \otimes u_p \otimes v_1 \otimes ... \otimes v_q \\
\end{array}$$
est une algèbre, appelée algèbre tensorielle de $V$.
\end{defi}

\begin{defi}[Algèbre extérieure]
Soit $V$ un $\R$-espace vectoriel et soit $k\in \N^*$, on appelle $k$-ème puissance extérieure de $V$, notée $\Lambda^k V$ le quotient de $\otimes^k V$ par l'espace vectoriel engendré par les éléments $v_1 \otimes ... \otimes v_k$ où il existe $i\neq j$ tels que $v_i = v_j$. Aussi on note $\Lambda^0 V = \R$.\\
On note $v_1 \wedge ... \wedge v_k$ la classe dans ce quotient de l'élément $v_1 \otimes ... \otimes v_k$. \\

De plus, si $V$ est de dimension finie et si $(e_1, ..., e_n)$ est une base de $V$ alors $e_{i_1} \wedge ... \wedge e_{i_k}$ pour tout $1 \leq i_1 < ... < i_k \leq n$ est une base de $\Lambda^k V$ qui est donc un espace vectoriel de dimension $\bino{k}{n}$. \\
On note $\Lambda V = \oplus_{k\geq 0} \Lambda^k V$. Le produit $\otimes : \otimes^p V\times \otimes^q V \to \otimes^{p+q} V$ passe au quotient en une application $\wedge : \Lambda^p V \times \Lambda^q V \to \Lambda^{p+q} V$ qui munit $\Lambda V$ d'une structure d'algèbre qu'on appelle algèbre extérieure de $V$.
\end{defi}

\begin{pro}
Soit $V$, $W$ et $F$ des $\R$-espaces vectoriels. On a les propriétés suivantes:
\be
	\item L'application $V^k \to \Lambda^k V, (v_1, ..., v_k) \mapsto v_1 \wedge ... \wedge v_k$ est $k$-linéaire alternée.
	\item (propriété universelle): $\Lambda^k V$ vérifie la proprété universelle: pour toute application $k$-linéaire alternée $f: V^k \to F$ il existe une unique application linéaire $g : \Lambda^k V \to F$ telle que le diagramme commute
	$$\begin{tikzcd}
	V^k \arrow[r, "f"]  \arrow[d, ".\wedge ."] & F \\
	\Lambda^kV \arrow[ru, "g"']& 
	\end{tikzcd}$$
	\item (fonctorialité):  Soit $f: V \to W$ une application linéaire, alors il existe une unique application linéaire $\Lambda^k f : \Lambda^k V \to \Lambda^k W$ définie par
	$$\Lambda^k f (v_1\wedge ... \wedge v_k) = f(v_1) \wedge ... \wedge f(v_k)$$
	et il existe une unique application linéaire $\Lambda f : \Lambda V \to \Lambda W$ définie par $\Lambda f = \oplus_{k\geq 0} \Lambda^k f$. \\
	\item (anticommutativité): Si $u \in \Lambda^p V$ et $v \in \Lambda^q V$ alors $u\wedge v = (-1)^{pq} v\wedge u$
	\item Si de plus $V$ est de dimension finie
	$$\Lambda^k V^* \simeq (\Lambda^k V)^*$$ 
	via l'isomorphisme $e_{i_1}^* \wedge ...\wedge e_{i_k}^* \mapsto (e_{i_1} \wedge ... \wedge e_{i_k})^*$ où $(e_i)_i$ est une base de $V$ et où $(e_i^*)_i$ désigne sa base duale dans $V^*$.
\ee
\end{pro}
\begin{preu}
(ii) Comme $f$ est $k$-linéaire par propriété universelle du produit tensoriel, il existe une unique $h: \otimes^k V \to F$ linéaire telle que $f(v_1, ..., v_k) = h(v_1 \otimes ... \otimes v_k)$. Puis comme $f$ est alternée, $h$ descend au quotient en une unique application $g: \Lambda^k V \to F$ linéaire telle que $f(v_1, ..., v_k) = g(v_1\wedge ...\wedge v_k)$.   
\end{preu}

De la même façon dont on identifie $E^*\otimes F^*$ avec les formes bilinéaires sur $E\times F$, on identifie $\otimes^k V^*$ avec les formes $k$-linéaires sur $V^k$ et on identifie $\Lambda^k V^*$ avec les formes $k$-linéaires alternées sur $V^k$. \\

De plus, on identifie également $\otimes^k V^* \otimes F$ avec les applications $k$-linéaire $V^k \to F$ et $\Lambda^k V^* \otimes F$ avec les application $k$-linéaire alternées $V^k \to F$. \\

On définit également le produit extérieur entre applications multilinéaires alternées en posant, 
$$\begin{array}{ccccc}
	\wedge & : & \Lambda^p V^* \otimes E \times \Lambda^q V^* \otimes F & \to & \Lambda^{p+q} V^* \otimes E \otimes F \\
		& & (\omega \otimes u, \eta \otimes v) & \mapsto & (\omega \wedge \eta) \otimes u \otimes v \\
\end{array}$$ \\

En particulier, on définit le produit extérieur entre applications multilinéaires alternées à valeur dans une algèbre de Lie $\mathfrak{g}$ en précomposant par le crochet de Lie. C'est à dire
$$\begin{array}{ccccc}
	\wedge & : & \Lambda^p V^* \otimes \mathfrak{g} \times \Lambda^q V^* \otimes \mathfrak{g} & \to & \Lambda^{p+q} V^* \otimes \mathfrak{g} \\
		& & (\omega \otimes u, \eta \otimes v) & \mapsto & (\omega \wedge \eta) \otimes [u, v] \\
\end{array}$$ \\

\section{Étoile de Hodge}

Soit $V$ un $\R$-espace vectoriel orienté de dimension $n$ muni d'une forme bilinéaire symétrique non dégénérée $\eta$. Notons $(e_1, ..., e_n)$ une base orthonormée directe de $V$. On pose 
$$\eta_k(u_1 \wedge ... \wedge u_k, v_1 \wedge ... \wedge v_k) = \det (\eta(u_i, v_j))$$
c'est aussi une forme bilinéaire symétrique non dégénérée sur $\Lambda^k V$.\\
Notons $\Vol = e_1 \wedge ... \wedge e_n$ appelé élément volume de $\Lambda^n V$. 

\begin{defi}[Étoile de Hodge]
Soit $V$ un $\R$-espace vectoriel orienté de dimension $n$ muni d'une forme bilinéaire symétrique non dégénérée $\eta$. Pour $0 \leq k \leq n$, il existe un unique isomorphisme linéaire $\star : \Lambda^k V \to \Lambda^{n-k} V$ appelée étoile de Hodge telle que 
$$\omega \wedge \rho = \eta_{n-k}(\star w, \rho) \Vol, \quad \forall \omega \in \Lambda^k V, \quad \forall \rho \in \Lambda^{n-k} V$$
\end{defi} 
\begin{preu}
Soit $\omega \in \Lambda^k V$ et on pose l'application linéaire $\rho \in \Lambda^{n-k} V \mapsto \omega \wedge \rho \in \Lambda^n V$. Comme $\Lambda^n V$ est de dimension 1, on a $\omega \wedge \rho = c(\rho)\Vol$. Alors $c : \Lambda^{n-k} V \to \R$ est une forme linéaire donc se représente par un unique $\star\omega \in \Lambda^{n-k} V$ via la forme bilinéaire symétrique non-dégénérée $\eta_{n-k}$.
\end{preu}

\begin{pro}
Soit $\omega = e_{i_1}\wedge ... \wedge e_{i_k} \in \Lambda^kV$ où $i_1 < ... < i_k$, alors 
$$\star \omega = \sigma e_{j_1} \wedge ... \wedge e_{j_{n-k}}$$
où $j_1 < ... < j_{n-k}$ complète la suite $i_1 < ... < i_k$ dans $\ieff{1}{n}$ et où le signe $\sigma$ est choisit tel que 
$$e_{i_1} \wedge ... \wedge e_{i_k} \wedge e_{j_1} \wedge ... \wedge e_{j_{n_k}} = \sigma \eta_{n-k}(e_{j_1} \wedge ... \wedge e_{j_{n_k}}, e_{j_1} \wedge ... \wedge e_{j_{n_k}})\Vol $$ 
\end{pro}